\ifdefined\iscombined
	\lecturetitle{Types of Testing}
\else
\documentclass{article}
\usepackage{listings}
\usepackage{amsthm}
\usepackage{comment}
\usepackage{amsmath}
\usepackage{amsfonts}
\usepackage{sectsty}
\usepackage{hyperref}
\usepackage{fancyhdr}
\usepackage[top=1in,bottom=1in,left=1in,right=1in]{geometry}
\usepackage{float}
\usepackage{graphicx}
\usepackage{tabularx}
\usepackage{mathtools}

\date{
	\vspace{-.75em}
	{\large \today}
}

\title{
	\vspace*{-1.5em}
	{\Huge Lecture 3} \\
	\vspace{-1em}
	\rule{\textwidth/4}{.01em} \\
	\vspace{-.25em}
	{\Large Types of Testing}
	\vspace{-.75em}
}

\author{
	{\large Matthew Farah}
}

\hypersetup{
	colorlinks=true,
	linkcolor=blue,
	filecolor=magenta,
	urlcolor=blue,
	citecolor=blue,
	anchorcolor=blue
}

\fancyhead{}
\fancyhead[L]{\today}
\fancyhead[R]{Lecture 3 - Types of Testing}

\sectionfont{\fontsize{12}{12}\bfseries}
\subsectionfont{\fontsize{10}{12}\normalfont}
\subsubsectionfont{\fontsize{10}{12}\normalfont}

\begin{document}

\maketitle
\pagestyle{fancy}

\fi

\begin{itemize}
	\item `Input cases' are groups of inputs to a program which cause it to behave similarly during execution.
	\begin{itemize}
		\item An `edge case' is an input belonging to either multiple input cases or none at all.
		\begin{itemize}
			\item Testing edge cases is key to ensuring programs behave as intended.
		\end{itemize}
		\item When input cases are well defined, a single representative test case from a group can be used to adequately verify that the program works as intended for all inputs in that test case's group.
	\end{itemize}
	\item Tests should ideally be written and execute as quickly as possible.
	\begin{itemize}
		\item However, testing may be constrained by its cost and the availability of code artifacts to test on.
	\end{itemize}
	\item The kinds of bugs being tested for dictate the kinds of tests to be used.
	\item Code `linting' is the act of enforcing certain standards and code style throughout a codebase.
	\item Certain kinds of tests require the program to be executed to work, while others do not.
	\begin{itemize}
		\item This splits tests into two distinct categories.
	\end{itemize}
	\item All tests are either:
	\begin{enumerate}
		\item Dynamic tests.
		\begin{itemize}
			\item Tests programs by executing them.
			\item Used to find runtime or integration bugs.
			\item Should be down throughout the Continuous Integration / Continuous Development (CI/CD) process.
		\end{itemize}
		\item Static tests.
		\begin{itemize}
			\item Tests programs \underline{without} executing them.
			\item Used to catch bugs early and perform code linting.
			\item Should be performed prior to any dynamic tests.
		\end{itemize}
	\end{enumerate}
	\item As well, certain kinds of tests work using or not using available knowledge about the code or its structure.
	\begin{itemize}
		\item This splits tests into another three distinct (somewhat more ambiguous) categories.
	\end{itemize}
	\item All tests are either:
	\begin{enumerate}
		\item Black-box tests.
		\begin{itemize}
			\item Tests programs without using \underline{any} knowledge of the code or its structure.
			\item Must be done as a dynamic test.
			\item Performed when:
			\begin{itemize}
				\item No knowledge of the code or its structure is available.
				\item The test is verifying that any changes made to the code or its stucture shouldn't affect its behaviour.
			\end{itemize}
		\end{itemize}
		\item Grey-box tests.
		\begin{itemize}
			\item Tests programs by using \underline{some} knowledge of the code and/or its structure.
			\item Typically done as a dynamic test.
			\item Performed when:
			%TODO: Add more
			\begin{itemize}
				\item The test is verifying that units are properly integrated.
			\end{itemize}
		\end{itemize}
		\item White-box tests.
		\begin{itemize}
			\item Tests programs by using \underline{full} knowledge of the code and its structure.
			\item Can be done as either a dynamic or static test.
			\item Performed when:
			%TODO: Add an example
			\begin{itemize}
				\item The test must do anything not possible by black-box or grey-box testing.
			\end{itemize}
		\end{itemize}
	\end{enumerate}
	\item Requirements define characteristics a given system must abide by.
	\item Requirements are divided into two distinct categories:
	\begin{enumerate}
		\item Functional requirements.
		\begin{itemize}
			\item Describe what the system must do.
			\begin{itemize}
				\item For example, the features or behaviour it must have.
			\end{itemize}
		\end{itemize}
		\item Non-functional requirements.
		\begin{itemize}
			\item A system's non-functional requirements describe the performance needs the system must satisfy.
			\begin{itemize}
				\item For example, the speed it must operate at.
			\end{itemize}
		\end{itemize}
	\end{enumerate}
	\item Tests can be written to verify that a system is compliant with its functional or non-functional requirements.
	\begin{itemize}
		\item This once again places tests into two distinct categories.
	\end{itemize}
	\item All tests are either:
	%TODO: Maybe talk about each a little more.
	\begin{enumerate}
		\item Functional tests.
		\item Non-functional tests.
	\end{enumerate}
\end{itemize}

\ifdefined\iscombined
	\newpage
\else
	\end{document}
\fi
